\chapter{A finer alphabet, and what it would buy}

\textbf{Purpose:} Work out what subdividing the octant alphabet would buy before anyone spends the compute on it, and record a derivation that did not hold along with what replaced it. \textbf{Scope:} the \(8 \times 4^k\) subdivision candidate, the rank bound, and 63 consecutive state pairs measured on both the working state and the finalized digest.

The eight-letter octant reading carries eight numbers about a source with \(256\) degrees of freedom. The obvious repair is a finer alphabet, and the candidate on the table is to subdivide each octant's spherical triangle, giving \(8 \times 4^k\) cells at depth \(k\).

\section{The answer}

Subdividing does not repair the reading, and the argument settling it is counting, and not statistics. The placement holds \(256\) points, and past a certain depth there are more cells than there are points to write into them.

\begin{center}
\begin{tabular}{@{}rrr@{}}
\toprule
\(k\) & cells & placement slots per cell \\
\midrule
0 & 8 & 32 \\
1 & 32 & 8 \\
2 & 128 & 2 \\
3 & 512 & 0.5 \\
\bottomrule
\end{tabular}
\end{center}

At \(k = 3\) there are twice as many cells as placement points, so most letters can never be written at all and the alphabet carries dead symbols. At \(k = 2\) the cell count equals the typical lit count, each cell holds two placement slots, and the reading has stopped being a coarse-graining: it is close to a permuted list of which bits are lit. On a 256-point placement \(k = 2\) is the last depth meaning anything, and by then the object being read is no longer a distribution over regions.

Raising the placement's point count moves the ceiling but not the problem, since the redraw level below is set by the ratio of cells to lit points and that ratio is what the subdivision changes.

\section{The same ceiling, counted on the reading instead of the placement}

The table above counts cells against placement points. The harmonic reading gives the same ceiling from the other side, and the two agree.

A reading to degree \(L\) carries exactly \((L+1)^2\) real numbers about its source. Against \(256\) degrees of freedom, degree eight is rank \(81\) and blind in \(175\) directions, degree fifteen is where \((L+1)^2\) first reaches \(256\), and the least singular value collapses to \(3.5 \times 10^{-3}\) there before recovering to \(1.5 \times 10^{-1}\) at sixteen. Refinement raises the count of numbers a reading carries; it does not raise the count past the source, and it makes the inversion worse on the way through the floor.

A finer alphabet is therefore bounded twice, once by how many cells the placement can fill and once by how many numbers the degree can carry, and neither bound is statistical.

\section{The scaling expectation}

Separately from the ceiling, a finer alphabet is expected to raise the delta and the redraw level together, leaving their ratio where it was.

The floor half of that is solid. Two independent states of the same weight give a level growing as the square root of the cell count, ordinary redraw statistics needing no assumption about the function being read. The delta half is only an expectation: a round's movement acts like a partial redraw, so it is expected to scale the same way. An attempt to derive that equality is recorded below, along with the reason it does not stand.

\section{A derivation that did not hold}

The argument was that the delta equals the floor times \(\sqrt{s}\), where \(s\) is the fraction of the lit set a round sends to fresh positions, on the grounds that points which did not move cancel cell by cell and leave the same calculation over \(s n\) points. The cell count then cancels between the two expressions and the ratio depends on \(s\) alone.

\textbf{The hole.} A share is a count divided by the weight, and the weight changes every round. On the working state it ranges from \(108\) to \(140\) lit bits and moves by \(4.3\) bits per round on average, \(3.5\%\) of the weight. A change in the denominator moves all eight shares even for points that never moved, so the cancellation the argument rests on does not happen.

\textbf{The test.} Predicting \(s\) from the ratio squared and counting the lit points that actually moved, over all \(63\) consecutive pairs, with the digest column kept beside it because the earlier version of this note was measured there:

\begin{center}
\begin{tabular}{@{}lll@{}}
\toprule
quantity & working state & finalized digest \\
\midrule
correlation of predicted against counted & \(0.105\) & \(0.068\) \\
worst predicted \(s\) & \(3.633\), an impossible value & \(3.508\) \\
observed ratio, per pair & \(0.382\) to \(1.906\) & \(0.507\) to \(1.873\) \\
\bottomrule
\end{tabular}
\end{center}

As a per-pair estimator it fails on either object, and it returns values above one where no fraction can go. Changing the object moved every figure and changed no conclusion. Running it on both objects earned that much.

\textbf{What holds.} On the mean, with both sides in the same units, the law's central claim stays close. The counted relocation in the law's own units averages \(1.0172\), whose square root is \(1.0085\), against an observed mean ratio of \(1.0361\). The two agree to \(2.7\%\).

The units matter and were the source of an earlier disagreement about this: a set difference between two states of weight \(n_2\) out of \(N\) saturates at \(1 - n_2/N\), about one half, and never at one, and a raw set difference of \(0.5\) already means complete independence. Comparing \(\sqrt{0.5} = 0.707\) against a ratio near \(1.0\) compares two different quantities.

The mean prediction therefore survives and the per-round estimator does not. Neither outcome supports using this argument as the reason a finer alphabet fails, and the two ceilings above are the reasons to give.

\section{The constant that replaced it}

The counted relocation is close to constant, the useful finding to come out of the test.

\begin{center}
\begin{tabular}{@{}p{0.34\linewidth}ll@{}}
\toprule
quantity & working state & finalized digest \\
\midrule
set difference between consecutive states & mean \(0.5215\), range \(0.421\) to \(0.619\) & mean \(0.4969\) \\
the same in matched units, where \(1.0\) is independence & \(1.0172\) & \(0.9976\) \\
bits differing between consecutive states & mean \(129.9\) of \(256\) & mean \(127.7\) \\
what a coin would give & \(128.0\) & \(128.0\) \\
\bottomrule
\end{tabular}
\end{center}

Every round relocates the lit set about as completely as an independent draw would, from the first round onward, with the working state sitting \(1.7\%\) above independence and the digest \(0.2\%\) below it. Reading \(s\) off the ratio is therefore estimating something that barely varies, and the backwards reading carries correspondingly little.

An earlier version of this note put worked examples in a table, \(s = 0.45\) at round 1 to 2 and \(s = 0.91\) across the first eight. Both were wrong, and they were wrong in a way worth naming: they came from reading the ratio off aggregate percentages instead of per pair, and they bracketed the true value by accident.

\section{What a frame actually is}

Worth stating, because it bounds every result above, and because an earlier version of this section had it backwards.

A frame is the eight working words as the round leaves them, with nothing added back. \texttt{rounds\_of} reads that by default. It can also return the finalized digest under \texttt{--read digest}, where \texttt{digest\_after} adds the initial values back into the working state before returning, and that feed-forward is a mixing step applied after the round.

The two objects are told apart by counting and not by assertion. A round only copies six of its eight words, \(b,c,d\) taking \(a,b,c\) and \(f,g,h\) taking \(e,f,g\). Over sixty-three transitions that is \(378\) carried words, and on the working state all \(378\) match the words they come from while on the digest none do.

That count also settles a question this section previously recorded as closed the wrong way. The register shift is present in the working state exactly, since matching all \(378\) carried words is what being a register shift consists of. The earlier text reported the shift as invisible, on a test that undid the word shift and found \(108.9\) of \(224\) bits differing against a coin's \(112.0\) with no exact match in sixty-four rounds. That test was run on digest frames, and it was measuring the feed-forward.

What is not established is a link between the shift and the octant delta. The delta's early excess over the redraw floor is measured, and a rigid relabeling of indices is the kind of move producing such an excess, but attributing the excess to the register shift is a reading of why and has not been measured. The two statements sit at different levels and only the first is a count.

\section{What would settle the subdivision}

One run at \(k = 1\) and \(k = 2\), each with its own redraw level computed at the matched weight, reporting the ratio and not the delta. Two ratios near the \(k = 0\) value confirm the expectation. A ratio falling with \(k\) refutes it cheaply, and the two ceilings still stand either way.

\textbf{Ownership.} The subdivision candidate and the run are the engine's. The refutation of the \(\sqrt{s}\) derivation is the engine's, and the rank bound is the engine's. The reproduction, the units correction, the object correction and the reruns above belong to this document.
